% Lösungen Einheit 3 von 4 – Flächen- und Volumeneinheiten (zusammenbau v0.1)
\einheitenkopf{Einheit 3 von 4 · Flächen- und Volumeneinheiten}

\erg{29}{a) hundert \quad b) $300$ dm² \quad c) $40\,000$ cm² \quad d) $700$ m² \quad e) mm³; cm³; dm³; m³ \quad f) $5\,000$ cm³ \quad g) $6$ dm³ \quad h) $1\,800$ ml \quad i) tausend \quad j) $1{,}25$ l $= 1\,250$ ml; $1\,250 : 200 = 6$ Rest $50$ – also $6$ volle Gläser \quad k) $75$ m³ $= 75\,000$ l; $75\,000 : 12\,500 = 6$ h}
\erg{30}{a) Fläche (Faktor zum Quadrat) \quad b) $5 \cdot 5 = 25$ cm², $6 \cdot 25 = 150$ cm² \quad c) $3 \cdot 2 + \frac{1}{2} \cdot 3 \cdot 1{,}5 = 8{,}25$ m²; $8{,}25 \cdot 4 = 33$ m³ \quad d) $g(x) - f(x) = -0{,}25x^4 + 0{,}75x^2 + 1$, Stammfunktion $-0{,}05x^5 + 0{,}25x^3 + x$; Wert bei $2$ minus Wert bei $-2$: $2{,}4 - (-2{,}4) = 4{,}8$ Flächeneinheiten; eine Flächeneinheit $= 30 \cdot 30 = 900$ cm², $4{,}8 \cdot 900 = 4\,320$ cm², beide Seiten $2 \cdot 4\,320 = 8\,640$ cm² $= 0{,}864$ m²}
\erg{31}{a) cm²; m²; a; ha; km²}
\erg{32}{a) $6$ cm$^2$; $18$ m$^2$; $7\,000$ m$^2$ (Briefmarke, Kinderzimmer, Fußballfeld)}
\erg{33}{a) Nils hat mit zehn statt mit hundert umgerechnet; Flächen haben die Umrechnungszahl hundert: $7 \cdot 100 = 700$, also $7$ dm² $= 700$ cm² \quad b) Ole hat den Maßstab nicht quadriert: eine Flächeneinheit sind $25 \cdot 25 = 625$ cm², also $8 \cdot 625 = 5\,000$ cm² \quad c) Ein Quadratmeter ist zehn Dezimeter lang und zehn Dezimeter breit: zehn Reihen mit je zehn Quadratdezimetern, also hundert. \quad d) Eine Flächeneinheit ist ein Quadrat mit zwei Seiten von je $3$ cm; Fläche ist Länge mal Länge, also $3$ cm mal $3$ cm. \quad e) $90 \cdot 40 \cdot 40 = 144\,000$ cm³ $= 144$ dm³ $= 144$ l \quad f) Tafel $3$ | $07$ | $15$, also $3{,}0715$ m² $= 30\,715$ cm²}

33f)

\sachtabelle{ccc}{m² & dm² & cm²}{3 & 07 & 15}

